\begin{abstract}
Positions of liquefied-natural-gas carriers are widely believed to carry
information about the price gap between North American and European gas: a cargo
that leaves the US Gulf today lands in Europe in two to three weeks, and that
supply is not yet in any published statistic. We test the belief directly. From
\nFixes{} AIS position fixes (\fixSpanLo{} to \fixSpanHi{}), reconstructed into
\nEvents{} port events and a \nSignalKeys-signal daily panel built twice --- once
with hindsight for validation and once strictly point-in-time for modelling --- we
ask whether tanker-derived signals predict the one- and four-week change in the
\HH{}--\TTF{} spread net of its own persistence, weather, storage and oil. Every
hypothesis, sign and acceptance bar was committed to a dated append-only log before
the corresponding fit. The answer is no. A controls-only model is worse than
assuming no change (\bMTwoSkillHOneCov{} at one week, \bMTwoSkillHFourCov{} at four);
no signal reaches a Newey--West $|t|$ of 1.96 (maximum \bMaxTCov); and three
pre-registered mechanisms that could hide an effect from a linear scan --- the
Europe-bound \emph{share} of cargo, an interaction with EU storage, and a
tightness-regime split --- fail under a Bonferroni bar with an untouched holdout.
The design detects effects of at least \powerFloor{} of residual variance and
observed at most \bMaxRsqCov, so the result rules out tradeable effects and is
silent on sub-tradeable ones. Alongside the null, we document two construction
hazards specific to voyage-derived features: the unit of observation is defined
retrospectively, and feature \emph{availability} changes across data vintages. A
naive loader leaked \leakFutureShare{} of the sample.
\end{abstract}
